% Additional exposition by Daniel Ortiz-Barrientos.
% Body only: requires amsmath, amssymb, natbib and hyperref in the wrapper.
% SF1 locators refer to Supplementary_File_1_Mathematics_Guide.tex,
% as prepared for the neutral preprint bundle on 15 November 2026.

A curved relationship between breeding value and expected reproduction can
have a close linear fit across a population. Understanding why requires us
to specify the quantity being fitted, the population weighting the fit, and
the prediction we want to make. This companion explains five choices behind
\emph{Additive Channels in Curved Fitness Landscapes}
\citep{OrtizBarrientosCooper2026}. Supplementary File~1 (SF1) supplies the
calculations; here we concentrate on why those calculations answer the
biological questions, and where a different question would require a
different approach.

\section{Which fitness quantity answers the question?}
\label{companion:target}

We want to describe how expected reproductive contributions differ among
individuals with different breeding values. Individuals sharing a breeding
value can differ in environment, phenotype and realised output. Averaging
over that variation defines the relationship of interest. With reproductive
output $W\geq0$ and breeding value $\mathbf A$, write
\begin{equation}
 W_A(\mathbf a)=\mathbb E[W\mid\mathbf A=\mathbf a],
 \qquad
 \ell_A(\mathbf a)=\log W_A(\mathbf a).
 \label{companion:eq:target}
\end{equation}
The logarithm requires a finite, positive conditional mean. The averaging includes
individuals that leave no offspring. This use of expected fitness at a
breeding value follows the quantitative-genetic treatment of environmental
variation \citep{Lande1979}.

The conditional mean preserves the expected reproductive contribution of
each breeding value. Weighting that value by its mean output gives the
same expected contribution as averaging the individual weighted outputs.
This follows from the law of conditional expectation. We can therefore
calculate parental moments without modelling every source of individual
variation in output. That variation still matters when estimating the
conditional mean from data.

Choosing breeding value identifies the genetic predictor whose relationship
with reproduction we wish to describe. It does not assume that all gene
action is additive. Breeding values and average effects refer to a
particular population \citep{Fisher1941}; treating their allelic coefficients
as fixed across generations will require an additional assumption in
Section~\ref{companion:inheritance}.

Why take a logarithm at all? The question concerns proportional differences
in expected output: a difference in $\ell_A$ is the logarithm of the ratio of
two conditional means. Multiplying every expected output by the same
positive constant therefore changes only the intercept. A quadratic log
mean also gives a positive expected output after exponentiation, consistent
with a log-link fitness model \citep{MorrisseyGoudie2022}. Modelling fitness
on its original scale is possible, but its slopes and explained variance
answer a different, absolute-scale question.

Taking logarithms of individual outputs first would change the target.
Generally, $\mathbb E[\log W\mid\mathbf A]$ differs from
$\log\mathbb E[W\mid\mathbf A]$. For two individuals leaving zero and two
offspring, the mean output is one and its logarithm is zero; the mean of
their individual logarithms is undefined. The required data model should
therefore describe conditional mean output, retaining zero counts, rather
than manufacture a logarithm for each count. SF1, Sections~2--3, develops
this distinction and the environmental averaging.

\clearpage
\section{Retain curvature, then assess the fit across the population}
\label{companion:projection}

A linear approximation cannot tell us when curvature matters, because it
has already omitted curvature. The next Taylor term is the smallest
extension that allows the slope to change. Centre breeding values at their
current mean $\mathbf a_0$, writing $\mathbf X=\mathbf A-\mathbf a_0$, and
form
\begin{equation}
 \widetilde\ell_A(\mathbf X)
 =\ell_0+\mathbf b^{\mathsf T}\mathbf X
       +\tfrac12\mathbf X^{\mathsf T}\mathbf H\mathbf X.
 \label{companion:eq:quadratic}
\end{equation}
Here $\mathbf b$ is the gradient and $\mathbf H$ the symmetric Hessian of
$\ell_A$ at the mean. With one trait these become the slope $b$ and
curvature $H$; bold symbols allow several traits. They describe the fitness
relationship, not allelic interactions. Twice differentiability justifies
this local expansion; it does not ensure that omitted terms are small
across the population. A more flexible fitted curve can assess that
approximation, at the cost of losing the following simple partition.

Why assess the fit with population least squares rather than the tangent
at the mean? A tangent describes change at one point. Our question concerns
variation among the breeding values actually present. Population least
squares minimises squared discrepancy under their frequencies, connecting
the calculation to regression descriptions of selection
\citep{LandeArnold1983,MooradWade2013}. This choice determines which
contributions count as explained and which remain after fitting the line.

The distribution of breeding values supplies those frequency weights. We
assume $\mathbf X\sim\mathcal N(\mathbf0,\mathbf G)$, so its covariance
$\mathbf G$ determines those weights. The Gaussian choice provides known
third and fourth moments; it is an assumption to assess, not a consequence
of defining breeding values. For a non-Gaussian population, additional
moments can change both the partition and the fitted slope.

We could instead calculate the best linear fit under an observed,
non-Gaussian distribution. That is a sensible choice when the immediate
question is how well a line describes that population. Here we also want to
ask what follows from changing covariance while keeping the fitness curve
fixed. A Gaussian family makes that comparison determinate: mean and
covariance specify the distribution. Outside that family, populations with
the same covariance can have different higher moments and different fits.

To separate variation from the mean, centre the quadratic contribution:
\begin{equation}
 L=\mathbf b^{\mathsf T}\mathbf X,
 \qquad
 Q_c=\tfrac12\left\{\mathbf X^{\mathsf T}\mathbf H\mathbf X
                         -\operatorname{tr}(\mathbf H\mathbf G)\right\}.
 \label{companion:eq:centred}
\end{equation}
The trace, $\operatorname{tr}$, adds a matrix's diagonal entries. With one
trait, the centred quadratic term is simply $Q_c=H(X^2-G)/2$, where $G$ is
the breeding-value variance. Subtracting its expectation changes the
intercept, not its variance.
Gaussian symmetry makes the third moments zero, so $L$ and $Q_c$ are
uncorrelated. The Gaussian fourth-moment identity \citep{Isserlis1918}
then gives
\begin{align}
 V_{\mathrm{lin}}&=\mathbf b^{\mathsf T}\mathbf G\mathbf b,
 & V_{\mathrm{quad}}&=\tfrac12\operatorname{tr}[(\mathbf H\mathbf G)^2],
 \label{companion:eq:components}\\
 \mathcal A_g&=\frac{V_{\mathrm{lin}}}
                    {V_{\mathrm{lin}}+V_{\mathrm{quad}}}.
 \label{companion:eq:share}
\end{align}
The denominator is the variance of the quadratic description. If both
components vanish, the share is undefined. Neither component is a variance
of realised offspring counts or a partition into additive and epistatic
genetic variance.

Because the centred quadratic term is uncorrelated with every linear
predictor, the best affine fit is
\begin{equation}
 \widehat{\widetilde\ell}_1(\mathbf X)
 =\ell_0+\tfrac12\operatorname{tr}(\mathbf H\mathbf G)
           +\mathbf b^{\mathsf T}\mathbf X.
 \label{companion:eq:affine}
\end{equation}
Its slope equals the local gradient, but its intercept includes the average
curvature contribution. Its explained variance is $V_{\mathrm{lin}}$, so
its population $R^2$ for the quadratic target is $\mathcal A_g$. This is
exact for that target under Gaussian weights; it is not automatically the
$R^2$ of the complete fitness curve or of a count-data model. SF1,
Section~4 and the multivariate treatment in Section~9, gives the intervening
moment calculations.

\section{Change population spread without changing the fitness curve}
\label{companion:scaling}

We can now isolate the population's contribution to the fit. Comparing
arbitrary populations would mix changes in their means, covariance
directions and fitness relationships. Instead, hold the mean, $\mathbf b$
and $\mathbf H$ fixed and multiply every variance and covariance by the
same positive number: $\mathbf G=s\mathbf G_0$. This controlled comparison
preserves the distribution's shape and orientation. Substitution gives
\begin{equation}
 V_{\mathrm{lin}}(s)=sV_{\mathrm{lin}}(1),
 \qquad
 V_{\mathrm{quad}}(s)=s^2V_{\mathrm{quad}}(1),
 \qquad
 \mathcal A_g(s)=\frac{V_{\mathrm{lin}}(1)}
 {V_{\mathrm{lin}}(1)+sV_{\mathrm{quad}}(1)}.
 \label{companion:eq:scaling}
\end{equation}
The linear contribution contains one factor of covariance and the
quadratic contribution two. When both are positive, reducing $s$ increases
the linear share even while the absolute linear variance falls.

For one trait, let $b=2$, $H=-0.5$ and initially $G=1$. Then
$V_{\mathrm{lin}}=4$, $V_{\mathrm{quad}}=0.125$ and
$\mathcal A_g=32/33$. Halving $G$ gives components $2$ and $0.03125$,
raising the share to $64/65$. The curve has not become less curved; the
population weights a narrower range of it.

This choice makes the cause of the change identifiable, but limits the
conclusion. Moving the mean can change the local derivatives, and rotating
the covariance can change the two components differently. Even a smaller
covariance in every direction need not reproduce proportional scaling.
At a peak where $\mathbf b=\mathbf0$, contraction alone does not create
linear variation. Nor must a biological change in spread leave the
conditional fitness relationship fixed. Reporting both variance components
alongside the share distinguishes these possibilities (SF1, Section~4).

Proportional scaling therefore answers a conditional question before any
evolutionary mechanism is proposed. It establishes what a change in spread
would do under specified conditions. Whether selection, demography or a
change in mating actually produces those conditions requires a separate
account. The comparison itself contains no inheritance model.

\section{Change the weights before predicting selected parents}
\label{companion:selection}

The fitted line was chosen to describe log expected fitness under population
frequencies. Selection asks how those frequencies change when individuals
contribute different amounts to reproduction. Expected output supplies the
additional weight:
\begin{equation}
 p_{\mathrm{sel}}(\mathbf x)
 =\frac{p(\mathbf x)W_A(\mathbf a_0+\mathbf x)}{\overline W},
 \qquad
 \overline W=\mathbb E[W].
 \label{companion:eq:weights}
\end{equation}
This is the reproductive weighting behind the covariance description of
selection \citep{Price1970}. It concerns parental breeding values represented
in reproductive contributions; transmission to offspring has yet to be
specified.

\emph{For the next calculation, assume that log expected fitness is exactly
quadratic over the Gaussian distribution.} A Taylor approximation near the
mean is insufficient: exponentiation can give substantial reproductive
weight to values poorly described by that approximation. For a more general
fitness curve, Equation~\eqref{companion:eq:weights} remains the starting
point, but its moments may require numerical integration.

Starting with a linear-only selection prediction would omit the effect we
want to assess. The exact quadratic model keeps curvature in the
reproductive calculation while leaving the selected distribution
analytically tractable. It therefore permits a direct comparison between
the two predictions without an additional weak-curvature approximation.
The cost is the stronger assumption about the whole fitness relationship.

\Needspace{11\baselineskip}
Work first with positive-definite $\mathbf G$, and write
$\boldsymbol\Gamma=-\mathbf H$. Multiplication of the Gaussian
density by the fitness weight adds their quadratic exponents. Their
combined coefficient is the precision
$\mathbf J=\mathbf G^{-1}+\boldsymbol\Gamma$, an inverse covariance.
It must be positive definite for the total weight to be finite. Completing
the square then identifies the selected moments:
\begin{equation}
 \mathbf G_{\mathrm{sel}}
    =(\mathbf G^{-1}+\boldsymbol\Gamma)^{-1},
 \qquad
 \boldsymbol\Delta_{\mathrm{full}}
    =(\mathbf G^{-1}+\boldsymbol\Gamma)^{-1}\mathbf b.
 \label{companion:eq:selected}
\end{equation}
These use established Gaussian, log-quadratic calculations
\citep{MorrisseyGoudie2022}. Under stabilising curvature
$\boldsymbol\Gamma\succeq0$, the positivity condition is automatic for
$\mathbf G\succ0$, and covariance decreases in every direction or stays
unchanged. The formulas also allow other curvature when $\mathbf J\succ0$,
but covariance need not decrease.

Return to $b=2$, $H=-0.5$, $G=1$. Omitting curvature predicts
$\Delta_{\mathrm{lin}}=Gb=2$, whereas
$\Delta_{\mathrm{full}}=4/3$. Thus
\begin{equation}
 \mathcal A_g=\frac{32}{33}\simeq0.97,
 \qquad
 \frac{|\Delta_{\mathrm{lin}}-\Delta_{\mathrm{full}}|}
      {|\Delta_{\mathrm{full}}|}=\frac12.
 \label{companion:eq:example}
\end{equation}
A large linear share has not controlled displacement error. Increasing
$|b|$ increases the share but leaves this relative error unchanged. The
variance comparison and the reproductive comparison use different weights
and different scales of fitness.

To assess displacement error in several traits, we must first choose its
scale. Ordinary Euclidean length would change if one trait were converted
from metres to centimetres. The population covariance supplies a biological
reference: measure a displacement relative to the variation available in
each combination of traits. The resulting norm is
$\|\mathbf v\|_{\mathbf G^{-1}}
=(\mathbf v^{\mathsf T}\mathbf G^{-1}\mathbf v)^{1/2}$.
Multiplying displacements by $\mathbf G^{-1/2}$ implements that choice and
makes the population covariance the identity. This is whitening; it changes
coordinates, not the population.

Define $\mathbf T=\mathbf G^{1/2}\boldsymbol\Gamma\mathbf G^{1/2}$ and
$\mathbf u_{\mathrm{full}}=\mathbf G^{-1/2}
\boldsymbol\Delta_{\mathrm{full}}$, with $\mathbf u_{\mathrm{lin}}$ defined
similarly. Equation~\eqref{companion:eq:selected} gives the useful relation
\begin{equation}
 \mathbf u_{\mathrm{lin}}-\mathbf u_{\mathrm{full}}
     =\mathbf T\mathbf u_{\mathrm{full}}.
 \label{companion:eq:error-relation}
\end{equation}
The spectral norm bounds the effect of $\mathbf T$ in any direction; its
Frobenius norm is the square root of the sum of squared entries
\citep{BoydVandenberghe2004}.
Since $\|\mathbf T\|_{\mathrm F}^2=2V_{\mathrm{quad}}$, a nonzero full
displacement satisfies
\begin{equation}
 \frac{\|\boldsymbol\Delta_{\mathrm{lin}}
             -\boldsymbol\Delta_{\mathrm{full}}\|_{\mathbf G^{-1}}}
      {\|\boldsymbol\Delta_{\mathrm{full}}\|_{\mathbf G^{-1}}}
 \leq\|\mathbf T\|_2
 \leq\sqrt{2V_{\mathrm{quad}}}.
 \label{companion:eq:bound}
\end{equation}
This bound uses absolute curvature variance. A small value guarantees a
small relative error; a large value need not be attained in the gradient's
direction. For singular $\mathbf G$, inverse operations are restricted to
directions with nonzero variance. If the full displacement is zero, relative
error is undefined. SF1, Sections~5 and~9, develops these selection results.

\clearpage
\section{Specify what reproduction transmits}
\label{companion:inheritance}

Selected-parent covariance cannot determine offspring covariance without
an inheritance model. Selection changes associations among alleles;
segregation and recombination can reduce those associations. The distinction
between genic variance and association variance is established in
quantitative genetics \citep{Bulmer1971,Bulmer1974,BartonTurelli1991}.
It suggests separating the genic baseline $\mathbf B$ from the departure
$\mathbf D=\mathbf G-\mathbf B$.

This separation determines what reproduction acts on in the model.
Multiplying the whole covariance by a retention fraction would also remove
part of the genic contribution. Retaining a fraction of its departure from
the baseline instead represents the reduction of associations. Reproduction
can then restore variance lost through those associations without
introducing new alleles or changing their effects.

We now assume fixed allelic effects, fixed $\mathbf B\succ0$ and fixed
stabilising curvature $\boldsymbol\Gamma\succeq0$. We also model
preselection breeding values as Gaussian in each generation, with exactly
quadratic log expected fitness. A rule for covariance alone does not
guarantee that offspring are Gaussian. Fixing effects does not fix genic
variance: under the random-mating diploid convention, a biallelic locus
contributes $2p(1-p)\alpha^2$, which changes with allele frequency $p$ even
if its effect $\alpha$ stays constant.
Holding $\mathbf B$ fixed isolates changes in associations. Following
allele frequencies explicitly would address additional biology, but would
require more state variables and a different recursion.

Assume that reproduction retains one fraction $0\leq\rho<1$ of the selected
departure from the baseline. Then
\begin{equation}
 \mathbf D_{t+1}=\rho(\mathbf G_{\mathrm{sel},t}-\mathbf B),
 \qquad
 \mathbf G_{t+1}
   =(1-\rho)\mathbf B
      +\rho(\mathbf G_t^{-1}+\boldsymbol\Gamma)^{-1}.
 \label{companion:eq:inheritance}
\end{equation}
This closure is an additional model assumption. One retention fraction
summarises the effects of locus-specific associations and reproduction;
it is not a universal conversion from recombination rate or mating system.
The familiar one-half benchmark belongs to the freely recombining,
random-mating infinitesimal setting.

We can iterate the update numerically, but a finite trajectory cannot
establish whether a positive equilibrium exists or all trajectories
converge. For matrices, compare variance in every direction using
$\mathbf C\preceq\mathbf D$ to mean that
$\mathbf D-\mathbf C$ is positive semidefinite. The continuous update maps
the closed, bounded, convex interval from $(1-\rho)\mathbf B$ to
$\mathbf B$ into itself. Reproduction supplies the lower bound; selection
cannot push a covariance already below $\mathbf B$ above it. This gives a
region in which to apply Brouwer's fixed-point theorem, establishing an
equilibrium \citep{GranasDugundji2003}; the covariance inequalities give
\begin{equation}
 (1-\rho)\mathbf B\preceq\mathbf G^*\preceq\mathbf B.
 \label{companion:eq:bounds}
\end{equation}
Existence and these bounds do not establish matrix uniqueness or
convergence.

For one trait, more can be shown without another long derivation. Writing
$g$ for variance and $\gamma\geq0$ for stabilising curvature, the update and
its derivative are
\begin{equation}
 f(g)=(1-\rho)B+\rho\frac{g}{1+\gamma g},
 \qquad
 0\leq f'(g)=\frac{\rho}{(1+\gamma g)^2}\leq\rho<1.
 \label{companion:eq:scalar}
\end{equation}
The separation of two trajectories after one update is at most $\rho$
times its previous value. The contraction theorem gives a unique
equilibrium and convergence
from every nonnegative starting variance \citep{GranasDugundji2003}.
In the recurring example, take $B=1$, $\gamma=0.5$ and $\rho=0.5$.
Selection changes variance from $1$ to $2/3$, then reproduction gives
$g_1=5/6$. Iteration approaches $g^*=(\sqrt{17}-1)/4\simeq0.781$.

\Needspace{6\baselineskip}
The boundary $\rho=1$ removes the positive baseline contribution and needs
separate treatment. Throughout, a covariance equilibrium is conditional on
the fixed parameters; it is not a joint equilibrium of mean, allele
frequencies and covariance. Calculating $\mathcal A_g$ through time also
requires the gradient history, which the covariance update does not supply.
SF1, Sections~6 and~9, gives the fixed-point arguments and boundary cases.

For estimation from observations, Supplementary File~3 works through
repeated measurements and offspring counts under a specified measurement
model, accounting for uncertainty in breeding values rather than treating
their estimates as error-free.

\paragraph{AI assistance.}
This additional exposition was developed through iterative discussion with
ChatGPT, guided by biological and mathematical questions set by Daniel
Ortiz-Barrientos. AI assistance supported drafting and explanation;
responsibility for the assumptions, arguments and interpretation remains
with the author.
